\documentclass[11pt,a4paper]{article}
\usepackage{fontspec,amsmath,amssymb,amsthm,microtype}
\setmainfont{lmroman10-regular.otf}[BoldFont=lmroman10-bold.otf,ItalicFont=lmroman10-italic.otf,BoldItalicFont=lmroman10-bolditalic.otf]
\usepackage[margin=25mm]{geometry}
\usepackage[unicode,hidelinks]{hyperref}
\hypersetup{pdfauthor={LangevinLean Project}}
\setlength{\emergencystretch}{3em}
\setlength{\parskip}{0.4em}
\setlength{\parindent}{0pt}
\allowdisplaybreaks[2]
\raggedbottom
\title{\texorpdfstring{Interacting particle Langevin: transport and entropy convergence}{Interacting particle Langevin: transport and entropy convergence}}
\author{LangevinLean Project}
\date{}
\begin{document}\maketitle
Draft mathematical exposition; these notes may contain mistakes.


The same Gaussian interacting-particle recursion admits two kinds of estimates. Joint strong convexity gives a quadratic-transport bound and parameter root mean square error about a marginal-likelihood maximizer. A logarithmic Sobolev inequality for the coupled particle target gives relative-entropy bounds for the joint and projected laws. The hypotheses, initial quantities, and particle normalizations differ and are stated separately.



\section{\texorpdfstring{The strongly convex transport and parameter estimates}{The strongly convex transport and parameter estimates}}


In scaled coordinates the summed potential has explicit curvature and gradient constants. These regularity properties normalize its Gibbs law and permit a stationary Langevin coupling for the full-state transport bound. Latent integration independently identifies the parameter Gibbs marginal and proves its concentration about the unique marginal-likelihood maximizer. Projection and Minkowski’s inequality then add equilibrium concentration to the transport estimate, giving the equilibrium, transient, and discretization terms in the parameter bound.




\section{\texorpdfstring{The joint and parameter entropy estimates}{The joint and parameter entropy estimates}}


The entropy estimate concerns the relative entropy of the current joint law with respect to the lifted Gibbs target. Applying data processing to both laws gives the projected parameter-law divergence. It does not by itself give concentration about a deterministic maximizing parameter. The conditional interpolation theorem and the Gaussian interpolation construction establish their respective entropy estimates under the distinct assumptions stated in each result.



\end{document}
